\chapter{Lezione 2 --- Spazi di funzioni, serie di Fourier, trasformate e distribuzioni}

\section{Spazio di funzioni periodiche}

Consideriamo uno spazio di funzioni periodiche di periodo \(T\).
L'idea è lavorare in uno spazio funzionale su cui sia possibile introdurre
una struttura geometrica analoga a quella degli spazi vettoriali.

Indichiamo con
\[
L_T^2
\]
lo spazio delle funzioni periodiche di periodo \(T\) che siano quadrato-sommabili
su un periodo, cioè tali che
\[
\int_{t_0}^{t_0+T} |f(t)|^2\,dt < +\infty.
\]

Poiché \(f\) è periodica, il valore di questo integrale non dipende dalla scelta di \(t_0\).

\section{Richiamo su integrabilità}

Negli appunti si distinguono due nozioni di integrale.

\subsection{Integrale di Riemann}

Nell'integrale di Riemann si approssima l'area sotto il grafico della funzione
discretizzando il dominio e costruendo somme di rettangoli.

\subsection{Integrale di Lebesgue}

Nell'integrale di Lebesgue la costruzione è più generale:
invece di discretizzare il dominio, si discretizza il codominio,
cioè si ragiona sui livelli assunti dalla funzione.

L'integrale di Lebesgue consente di trattare classi più ampie di funzioni
e, sotto ipotesi opportune, permette di scambiare limiti e integrali.

Nel seguito assumiamo che l'integrabilità sia nel senso di Lebesgue.

\section{Perché funzioni a quadrato sommabile}

Vogliamo introdurre nello spazio delle funzioni una struttura geometrica,
in particolare prodotto scalare e norma.

Per questo consideriamo funzioni
\[
f:\mathbb{R}\to\mathbb{R}
\quad \text{oppure} \quad
f:\mathbb{R}\to\mathbb{C}
\]
tali che
\[
\int_{t_0}^{t_0+T} |f(t)|^2\,dt < +\infty.
\]

Questa ipotesi garantisce che la norma della funzione sia finita.

\section{Prodotto scalare in \(L_T^2\)}

Se \(f_1,f_2\in L_T^2\), definiamo il prodotto scalare come
\[
\langle f_1,f_2\rangle
=
\frac{1}{T}\int_{t_0}^{t_0+T} f_1(t)\,\overline{f_2(t)}\,dt.
\]

Se le funzioni sono reali, la coniugazione si può omettere:
\[
\langle f_1,f_2\rangle
=
\frac{1}{T}\int_{t_0}^{t_0+T} f_1(t)\,f_2(t)\,dt.
\]

Il fattore \(1/T\) serve a normalizzare il prodotto scalare rispetto alla lunghezza del periodo.

\section{Norma}

Una volta definito il prodotto scalare, la norma associata è
\[
\|f\|=\sqrt{\langle f,f\rangle}.
\]

Questo chiarisce perché si richieda che le funzioni appartengano a \(L_T^2\):
senza quadrato-sommabilità la norma potrebbe non esistere oppure essere infinita.

\section{Interpretazione geometrica}

Con prodotto scalare e norma, lo spazio delle funzioni acquista una struttura
simile a quella di uno spazio vettoriale euclideo.

Di conseguenza si può parlare di:
\begin{itemize}
    \item basi;
    \item ortogonalità;
    \item ortonormalità;
    \item proiezioni.
\end{itemize}

L'idea fondamentale di Fourier è che una funzione periodica possa essere espressa
come combinazione lineare di sinusoidi a frequenze diverse.

\section{Formula di Eulero}

Ricordiamo la formula di Eulero:
\[
e^{j\omega t}=\cos(\omega t)+j\sin(\omega t).
\]

Questa scrittura permette di trattare seni e coseni con un'unica notazione
tramite esponenziali complessi.

\section{Serie di Fourier complessa}

Sia \(f\in L_T^2\). Allora, in forma complessa, la sua serie di Fourier è
\[
f(t)=\sum_{n=-\infty}^{+\infty} C_n e^{j n\omega_0 t},
\qquad
\omega_0=\frac{2\pi}{T}.
\]

Qui:
\begin{itemize}
    \item \(T\) è il periodo della funzione;
    \item \(\omega_0\) è la pulsazione fondamentale;
    \item la frequenza fondamentale in Hertz è
    \[
    \nu_0=\frac{1}{T}=\frac{\omega_0}{2\pi}.
    \]
\end{itemize}

La formula dice che ogni funzione del nostro spazio si rappresenta come combinazione lineare
di armoniche della fondamentale.

\section{Che funzioni base compaiono}

Le funzioni base sono
\[
e^{j n\omega_0 t},
\qquad n\in\mathbb{Z}.
\]

Equivalentemente si può parlare di seni e coseni alle pulsazioni
\[
n\omega_0.
\]

Queste sono precisamente le sinusoidi che hanno periodo \(T\) oppure un suo sottomultiplo:
\[
\frac{T}{|n|}, \qquad n\neq 0.
\]

La famiglia delle funzioni base è infinita ma numerabile,
perché è indicizzata dagli interi.

\section{Ortogonalità della base}

Nel prodotto scalare definito sopra, le funzioni esponenziali complesse risultano ortogonali:
\[
\left\langle e^{j m\omega_0 t}, e^{j n\omega_0 t}\right\rangle
=
\frac{1}{T}\int_{t_0}^{t_0+T}
e^{j m\omega_0 t}\,\overline{e^{j n\omega_0 t}}\,dt.
\]

Poiché
\[
\overline{e^{j n\omega_0 t}}=e^{-j n\omega_0 t},
\]
si ottiene
\[
\left\langle e^{j m\omega_0 t}, e^{j n\omega_0 t}\right\rangle
=
\frac{1}{T}\int_{t_0}^{t_0+T} e^{j(m-n)\omega_0 t}\,dt
=
\begin{cases}
1, & m=n,\\
0, & m\neq n.
\end{cases}
\]

Quindi la famiglia
\[
\left\{e^{j n\omega_0 t}\right\}_{n\in\mathbb{Z}}
\]
è ortonormale.

\section{Coefficienti della serie di Fourier}

I coefficienti \(C_n\) si ottengono proiettando \(f\) sulla funzione base corrispondente:
\[
C_n=\left\langle f(t), e^{j n\omega_0 t}\right\rangle.
\]

Usando la definizione di prodotto scalare si ricava
\[
C_n
=
\frac{1}{T}\int_{t_0}^{t_0+T} f(t)e^{-j n\omega_0 t}\,dt.
\]

Scegliendo per comodità \(t_0=0\), si ha
\[
C_n
=
\frac{1}{T}\int_0^T f(t)e^{-j n\omega_0 t}\,dt.
\]

\section{Significato dei coefficienti \(C_n\)}

In generale \(C_n\) è un numero complesso.

Possiamo quindi scrivere
\[
C_n = |C_n| e^{j\varphi_n}.
\]

Il modulo \(|C_n|\) misura il peso dell'armonica \(n\omega_0\),
mentre la fase \(\varphi_n\) ne misura lo sfasamento.

Se una funzione coincide con una singola armonica, allora un solo coefficiente è non nullo;
negli altri casi il segnale è una combinazione lineare di più armoniche.

\section{Interpretazione per i sistemi lineari}

Se un segnale periodico entra in un sistema lineare, possiamo scomporlo in armoniche.
Per linearità, la risposta complessiva si ottiene sommando le risposte del sistema
alle singole componenti sinusoidali.

Questa è la base concettuale dell'analisi in frequenza dei sistemi lineari.

\section{Dal caso periodico al caso non periodico}

Se il periodo \(T\) cresce, la pulsazione fondamentale
\[
\omega_0=\frac{2\pi}{T}
\]
diventa sempre più piccola.

Nel limite \(T\to+\infty\):
\begin{itemize}
    \item le frequenze discrete \(n\omega_0\) si addensano;
    \item la sommatoria sulle armoniche tende a un integrale;
    \item si passa da una descrizione discreta in frequenza a una descrizione continua.
\end{itemize}

Questo passaggio porta dalla serie di Fourier alla trasformata di Fourier.

\section{Trasformata di Fourier}

Nel caso non periodico si scrive formalmente
\[
f(t)=\frac{1}{2\pi}\int_{-\infty}^{+\infty} F(\omega)e^{j\omega t}\,d\omega,
\]
dove
\[
\mathcal{F}\{f(t)\}=F(\omega)
=
\int_{-\infty}^{+\infty} f(t)e^{-j\omega t}\,dt.
\]

Qui \(\omega\) è una variabile continua e \(F(\omega)\) è, in generale, una funzione complessa.

Possiamo scrivere
\[
F(\omega)=|F(\omega)|e^{j\arg F(\omega)}.
\]

Il modulo \(|F(\omega)|\) descrive il contenuto in frequenza del segnale;
la fase \(\arg F(\omega)\) descrive lo sfasamento delle varie componenti.

\section{Segnali stazionari e significato della Fourier}

La trasformata di Fourier è lo strumento naturale per analizzare segnali stazionari
oppure segnali definiti su tutto l'asse reale.

Nel caso di segnali periodici o stazionari, il contenuto in frequenza non dipende
dall'istante da cui si comincia a osservare il segnale.

Per questo la Fourier è particolarmente adatta a descrivere il regime
e il contenuto armonico del segnale.

\section{Trasformata di Laplace}

La trasformata di Laplace unilaterale si definisce come
\[
\mathcal{L}\{f(t)\}=F(s)
=
\int_{0^-}^{+\infty} f(t)e^{-st}\,dt,
\]
dove
\[
s\in\mathbb{C},
\qquad
s=\sigma+j\omega.
\]

La variabile complessa \(s\) permette di pesare il segnale non solo con sinusoidi,
ma con esponenziali del tipo \(e^{st}\).

\section{Significato della Laplace}

La trasformata di Laplace è particolarmente utile per:
\begin{itemize}
    \item segnali causali;
    \item fenomeni transitori;
    \item studio dei sistemi dinamici;
    \item trattamento delle condizioni iniziali;
    \item descrizione di segnali che non ammettono Fourier classica.
\end{itemize}

Per questo, dal punto di vista fisico e ingegneristico:
\begin{itemize}
    \item Fourier è soprattutto uno strumento per il contenuto in frequenza
    e per i segnali stazionari;
    \item Laplace è soprattutto uno strumento per sistemi causali, transitori
    e risposte forzate.
\end{itemize}

\section{Rapporto tra Fourier e Laplace}

Formalmente, la trasformata di Fourier si può vedere come la trasformata di Laplace
valutata sull'asse immaginario:
\[
s=j\omega.
\]

Tuttavia questo passaggio è lecito solo quando l'asse immaginario appartiene
alla regione di convergenza della trasformata di Laplace.

Quindi Fourier non è semplicemente ``Laplace con \(s=j\omega\)''
in modo automatico: il legame è corretto, ma richiede condizioni di esistenza.

\section{Distribuzioni}

Per trattare in modo rigoroso oggetti come impulso, gradino e loro derivate,
si introducono le distribuzioni.

Una distribuzione non è, in generale, una funzione regolare nel senso classico,
ma un oggetto definito attraverso il suo effetto integrale su funzioni test.

\section{Delta di Dirac}

L'esempio fondamentale è la delta di Dirac \(\delta(t)\), definita dalla proprietà
\[
\int_{-\infty}^{+\infty} f(t)\delta(t)\,dt = f(0).
\]

Questa relazione dice che la delta estrae il valore della funzione nel punto \(t=0\).

Intuitivamente:
\begin{itemize}
    \item la delta è concentrata in \(t=0\);
    \item il suo effetto è nullo ovunque tranne che in quell'istante;
    \item ha senso solo sotto segno di integrale.
\end{itemize}

\section{Altre distribuzioni utili}

Anche il gradino, la rampa e le derivate della delta
si interpretano correttamente nel quadro delle distribuzioni.

Questo è importante perché trasformate come Laplace e Fourier,
essendo operazioni integrali, possono essere estese anche a questi oggetti generalizzati.

\section{Sintesi della lezione}

I punti chiave della lezione sono:
\begin{enumerate}
    \item introduzione dello spazio \(L_T^2\) delle funzioni periodiche quadrato-sommabili;
    \item definizione di prodotto scalare e norma;
    \item interpretazione geometrica dello spazio delle funzioni;
    \item base ortonormale di esponenziali complessi;
    \item serie di Fourier e coefficienti \(C_n\);
    \item passaggio dal caso periodico al caso non periodico;
    \item trasformata di Fourier;
    \item trasformata di Laplace;
    \item distribuzioni e delta di Dirac.
\end{enumerate}